% =====================================================================
%  Capitulo 3 --- Planning
%  Transcricao fiel de src/04-planning.html (secoes 3.1 a 3.3)
% =====================================================================
\chapter{Planning}

\section{Values, mission and vision}
AstraZeneca does not use the terms ``mission'' and ``vision''. It uses a \textit{Purpose}, which plays a role similar to the mission by explaining why the organisation exists, the \textit{Ambition 2030}, which corresponds to the vision and sets the intended direction, and five \textit{Values}, which guide the organisational culture and decision-making \cite[pp.~8, 10]{AR}, \cite[pp.~15--16]{PLAN}.

\subsection{Mission}
``We push the boundaries of science to deliver life-changing medicines'' \cite[p.~8]{AR}. The statement is short and clear about why the company exists: to create medicines that change patients' lives. The Group also states that it aims to create value for patients, society, the planet and shareholders \cite[p.~8]{AR}, which shows that it recognises several stakeholders and not only its shareholders. The scope is not part of the mission, but it can be completed with the business model: research, development, manufacturing and commercialisation of prescription medicines across primary care, specialty care and rare diseases \cite[p.~8]{AR}.

\subsection{Vision}
To be pioneers in science, lead in its disease areas and transform patient outcomes, with concrete targets for 2030: at least 20 new medicines and \$80bn in Total Revenue \cite[p.~10]{AR}. The numerical targets make the vision measurable and allow progress to be monitored. These targets are used as the strategic objective in Section~\ref{sec:objetivos}. The overall vision is then broken down into ambitions for each therapy area:

\begin{itemize}
  \item in Oncology, ``eliminate cancer as a cause of death'' \cite[p.~14]{AR};
  \item in BioPharmaceuticals, ``transform care for billions of people living with chronic diseases'' \cite[p.~16]{AR};
  \item in Rare Disease, ``pioneering new possibilities'', focusing on diseases with significant unmet medical need where scientific progress has been absent or limited \cite[p.~22]{AR}.
\end{itemize}

Each unit therefore receives its own direction, aligned with the Group's vision.

\subsection{Values}

\begin{xltabular}{\linewidth}{@{}lY@{}}
\caption{AstraZeneca's values. Source: \cite[p.~8]{AR}.}\label{tab:valores}\\*
\toprule
\textbf{Value} & \textbf{Meaning} \\*
\midrule
\endfirsthead
\toprule
\textbf{Value} & \textbf{Meaning} \\*
\midrule
\endhead
\bottomrule
\endfoot
We follow the science & Push the boundaries of science and work creatively with partners \\
We put patients first & Understand patients' needs and consider them in every decision \\
We play to win & Build high-performing, inclusive and diverse teams and make the right choices to win \\
We do the right thing & Apply high ethical standards in all aspects of the business \\
We are entrepreneurial & Act with urgency, bravery and resilience, taking smart risks \\
\end{xltabular}

The focus on science and patients is directly linked to the research and development of medicines, while integrity and entrepreneurship support responsible conduct and the search for new opportunities. The values guide the company's culture and decisions and are measured in the internal Pulse survey: 88\% of employees see the company as patient-oriented \cite[p.~38]{AR}. Measuring them is a form of control, because it allows the company to check whether the culture matches what is intended.

\begin{analysis}
Mission, vision and values are consistent with each other and share the same central idea: science in the service of the patient. The main challenge for management is to reconcile \textit{play to win} with \textit{put patients first}, that is, to balance the profitability required by shareholders with patients' access to medicines, in a context of pressure on prices \cite[pp.~32--33]{AR}.
\end{analysis}

\section{Objectives at the three management levels}\label{sec:objetivos}
Strategic objectives set the overall long-term results, tactical objectives translate the strategy into departmental plans, and operational objectives guide concrete day-to-day activities \cite[pp.~5--9]{PLAN}. AstraZeneca discloses only its strategic objectives. The tactical and operational objectives below are proposed here, not disclosed by the company, designed so that each one contributes to the level above. The Ambition 2030 is pursued through the \textit{Growth Through Innovation} strategy, which has three priorities: Science and Innovation, Growth and Therapy Area Leadership, and People and Sustainability \cite[p.~10]{AR}. The tactical objective below falls under the second priority and the operational objective supports the first.

\begin{xltabular}{\linewidth}{@{}N{3.1cm}Yl@{}}
\caption{One objective for each management level. Only the strategic objective is declared by the Group.}\label{tab:objetivos}\\*
\toprule
\textbf{Level} & \textbf{Objective} & \textbf{Horizon} \\*
\midrule
\endfirsthead
\toprule
\textbf{Level} & \textbf{Objective} & \textbf{Horizon} \\*
\midrule
\endhead
\bottomrule
\endfoot
\textbf{Strategic}\newline CEO and SET & Launch at least 20 new medicines and reach \$80bn in Total Revenue by 2030 \cite[p.~10]{AR} & Long \\
\textbf{Tactical}\newline Head of oncology clinical development & Submit at least three new oncology medicines for regulatory approval by the end of 2027 & Medium \\
\textbf{Operational}\newline Team leader of a clinical trial & Review at least 98\% of the trial's documents within the deadline, with no critical findings, by the end of the next quarter & Short \\
\end{xltabular}

\textbf{Strategic.} It covers the whole company, has a long horizon and requires decisions only top management can take, such as investment in R\&D, which areas to prioritise and acquisitions.

\textbf{Tactical.} It turns the strategic target into a result for one specific unit, Oncology, which is the largest area (44\% of revenue) \cite[p.~2]{AR} and had 10 positive Phase III readouts in 2025 \cite[p.~13]{AR}. It belongs to a middle manager, who coordinates teams but does not set the Group's strategy.

\textbf{Operational.} It focuses on a concrete task, has a short horizon and a precise indicator. The quality of documentation is essential for regulatory approval, and the Group already reports zero critical observations in 42 external inspections \cite[p.~4]{AR}.

\begin{analysis}
The three objectives are SMART and form a cascade. Reliable documentation (operational) enables regulatory submissions (tactical), and approved submissions become new medicines and revenue (strategic).
\end{analysis}

\section{SWOT analysis and strategy}
The SWOT is built by looking outside the company first and then inside it, following the course material \cite[pp.~17--29]{PLAN}. Outside, PESTEL covers the wider environment \cite[pp.~19--21]{PLAN} and Porter's five forces cover the pharmaceutical sector \cite[pp.~21--24]{PLAN}. Inside, the value chain shows what the company does \cite[p.~25]{PLAN} and VRIO shows which resources give it an edge \cite[p.~26]{PLAN}. Strengths and weaknesses describe the company today; opportunities and threats describe what is coming \cite[p.~27]{PLAN}. The facts come from the Annual Report and other cited sources, while ratings and classifications are an assessment, not company statements.

\begin{figure}[htbp]\centering
\begin{tikzpicture}[node distance=7mm and 6mm]
\node[fl] (a1) {PESTEL\\{\color{azmuted}wider environment}};
\node[fl, right=of a1] (a2) {Five forces\\{\color{azmuted}the sector}};
\node[fl, right=of a2] (a3) {Opportunities \& Threats\\{\color{azmuted}outside, future}};
\node[fl, below=of a1] (b1) {Value chain\\{\color{azmuted}activities}};
\node[fl, right=of b1] (b2) {VRIO\\{\color{azmuted}resources}};
\node[fl, right=of b2] (b3) {Strengths \& Weaknesses\\{\color{azmuted}inside, today}};
\node[fle, right=of a3] (sw) {SWOT};
\node[fle, right=of b3] (st) {Crossing\\$\rightarrow$ Strategy};
\draw[se] (a1)--(a2); \draw[se] (a2)--(a3); \draw[se] (a3)--(sw);
\draw[se] (b1)--(b2); \draw[se] (b2)--(b3); \draw[se] (b3)--(st);
\draw[se] (sw)--(st);
\end{tikzpicture}
\caption{From the environment to the strategy, following the method in the course material \cite[pp.~17--29]{PLAN}.}
\label{fig:metodo}
\end{figure}
\subsection{The wider environment: PESTEL}
The picture outside AstraZeneca is mixed. Demand is on the company's side: people live longer, chronic diseases and cancer keep growing, and health spending keeps rising. Price and regulation work against it: health budgets are tight, the US now negotiates prices, and governments want to limit patent protection. The course material asks us to rate the impact of each factor and to look at how factors interact \cite[p.~19]{PLAN}; Table~\ref{tab:pestel} sets this out factor by factor.


\begin{xltabular}{\linewidth}{@{}N{2.6cm}Yp{2.1cm}p{1.5cm}@{}}
\caption{PESTEL analysis. Inflation, education levels and government stability are left out because the Annual Report does not treat them as business factors.}\label{tab:pestel}\\*
\toprule
\textbf{Factor} & \textbf{Evidence} & \textbf{Class} & \textbf{Impact} \\*
\midrule
\endfirsthead
\toprule
\textbf{Factor} & \textbf{Evidence} & \textbf{Class} & \textbf{Impact} \\*
\midrule
\endhead
\bottomrule
\endfoot
\multicolumn{4}{@{}l@{}}{\textbf{Political}} \\*
\midrule
Geopolitics & Tension and strategic rivalry between the US and China, and rising economic nationalism \cite[p.~6]{AR}. Trade controls, tariffs and taxes may raise costs or reduce revenue \cite[p.~49]{AR}. & \bbad{Threat} & High \\
Localisation and tariffs & Governments encourage companies to localise operations to reduce supply-chain risk \cite[p.~7]{AR}. The Group obtained a three-year exemption from Section 232 tariffs and plans to invest \$50bn in US manufacturing and R\&D \cite[pp.~4, 32--33]{AR}. & \bgood{Opportunity} & Medium \\
Taxation & The Group paid \$2,845m of income tax in 2025. The Core tax rate was 18\% and is guided to 18--22\% in 2026 \cite[pp.~59, 64]{AR}. & \bneut{Neutral} & Low \\
\midrule
\multicolumn{4}{@{}l@{}}{\textbf{Economic}} \\*
\midrule
Budgets and pricing & Healthcare budgets are under pressure, leading to downward pressure on pricing \cite[p.~7]{AR}. Payer consolidation and restrictive reimbursement drive significant pricing pressure \cite[p.~32]{AR}. & \bbad{Threat} & High \\
Market growth & Global healthcare spending is expected to grow 7.1\% a year in 2025--2029, and global pharmaceutical sales grew 10.0\% in 2025 \cite[p.~6]{AR}. Global GDP growth is forecast at 3.3\% in 2026 and 3.2\% in 2027 \cite[p.~7]{AR}. & \bgood{Opportunity} & High \\
Shift to emerging markets & Economic power is shifting from the G7 to large emerging markets such as China, India and Indonesia \cite[p.~6]{AR}. & \bgood{Opportunity} & Medium \\
Interest rates and capital & Borrowing under committed facilities is priced on SOFR plus a margin, and treasury manages interest-rate, currency and credit risk \cite[pp.~61, 64]{AR}. The Group holds \$4,875m of undrawn committed facilities with no covenants \cite[p.~61]{AR}. & \bneut{Neutral} & Low \\
Currency & Significant changes in exchange rates may adversely affect the business \cite[p.~49]{AR}. At January 2026 rates, 2026 Total Revenue would gain a low single-digit percentage \cite[p.~64]{AR}. & \bneut{Neutral} & Low \\
\midrule
\multicolumn{4}{@{}l@{}}{\textbf{Social}} \\*
\midrule
Ageing & The population aged 65 and over is set to triple by 2100, which raises demand for care \cite[pp.~6--7]{AR}. & \bgood{Opportunity} & High \\
Disease burden & Non-communicable diseases are 75\% of non-pandemic deaths. About 1.4 billion people have hypertension, over 30 million new cancer diagnoses are estimated by 2040, and fewer than 10\% of rare diseases have an approved treatment \cite[pp.~7, 13, 17, 23]{AR}. & \bgood{Opportunity} & High \\
Talent & The Group builds talent internally: 5,100 promotions and 2,760 external hires in its global hubs in 2025 \cite[p.~39]{AR}. & \bneut{Neutral} & Low \\
\midrule
\multicolumn{4}{@{}l@{}}{\textbf{Technological}} \\*
\midrule
Artificial intelligence & Generative AI is estimated to create \$4--7bn of value a year for pharma, and the AR notes that AI ``will unlock value and bring new risks'' \cite[p.~7]{AR}. & Opportunity with risk & Medium \\
R\&D capacity and collaboration & R\&D expense was \$14,232m and more than 100 Phase III studies are ongoing \cite[pp.~3, 55]{AR}. More than 1,000 collaborations give access to technology \cite[p.~37]{AR}. & \bgood{Opportunity} & Medium \\
\midrule
\multicolumn{4}{@{}l@{}}{\textbf{Environmental}} \\*
\midrule
Climate & The climate crisis worsens health inequities and strains health systems \cite[p.~7]{AR}. The Group has its own science-based targets, including net zero by 2045 \cite[p.~42]{AR}. Extreme weather threatens sites, suppliers and distribution routes \cite[p.~44]{AR}. & \bbad{Threat} & Medium \\
Water, waste and nature & Operations cut water use by 24\% and waste by 23\% against 2015 \cite[p.~33]{AR}. The Group recognises the active ingredients in its products as its most material water-pollution impact \cite[p.~45]{AR}. & \bbad{Threat} & Low \\
\midrule
\multicolumn{4}{@{}l@{}}{\textbf{Legal}} \\*
\midrule
Price regulation (IRA) & In the US, the Inflation Reduction Act subjects medicines to Medicare-negotiated prices. Farxiga was selected and Calquence takes effect in 2027. The AR expects the impact on Farxiga to be manageable \cite[p.~33]{AR}; it is rated high here, because Farxiga is the largest product at \$8,492m, 14.5\% of Total Revenue \cite[p.~56]{AR}. & \bbad{Threat} & High \\
Intellectual property & Governments and payers press to limit intellectual-property protection \cite[p.~49]{AR}. & \bbad{Threat} & High \\
EU pharmaceutical reform & The December 2025 provisional agreement sets eight years of data protection plus one year of market protection (extendable by one year if two of three conditions are met), and a Bolar exemption so that generics can be available on day one after intellectual-property expiry \cite{EU}. & \bbad{Threat} & Medium \\
Sector regulation & Delays in regulatory approvals could delay marketing and reduce revenue \cite[p.~48]{AR}. Products need approval from authorities such as the FDA and the EMA \cite[p.~20]{AR}. & \bbad{Threat} & Medium \\
Data and AI rules & A growing range of legislative and regulatory requirements covers data and AI \cite[p.~48]{AR}. & \bbad{Threat} & Low \\
\end{xltabular}

\begin{analysis}
The factors do not act alone. US tariffs, the pricing agreement and local manufacturing come as one package: a tariff exemption in exchange for prices aligned with other wealthy countries and medicines made in the US \cite[pp.~4, 32--33]{AR}. In 2026 the IRA price for Farxiga arrives in the same year it loses US exclusivity, although the AR expects the impact to be manageable \cite[p.~33]{AR}. Ageing pushes demand up while tight budgets hold prices down, which is why the Group uses tiered pricing \cite[pp.~6--7, 41]{AR}. Looking ahead, health spending keeps growing to 2029, Calquence enters price negotiation in 2027 and the net zero target runs to 2045 \cite[pp.~6, 33, 42]{AR}.
\end{analysis}

\subsection{The sector: Porter's five forces}
This section looks at the market for innovative prescription medicines. Each force is rated strong, moderate or weak \cite[p.~21]{PLAN}, using the factors that the course material lists for each one \cite[pp.~22--24]{PLAN}.


\begin{xltabular}{\linewidth}{@{}N{2.5cm}Yp{2.3cm}@{}}
\caption{Porter's five forces. Intensity is not rated in the Annual Report.}\label{tab:porter}\\*
\toprule
\textbf{Force} & \textbf{Evidence} & \textbf{Intensity} \\*
\midrule
\endfirsthead
\toprule
\textbf{Force} & \textbf{Evidence} & \textbf{Intensity} \\*
\midrule
\endhead
\bottomrule
\endfoot
\textbf{Rivalry} & \textit{Number and balance:} several large groups of similar size. In 2024 the top 10 earned between \$44.47bn (Sanofi) and \$64.17bn (Merck \& Co.), and AstraZeneca was 5th with \$54.07bn, behind Merck \& Co., Pfizer, J\&J and AbbVie \cite{P50}. \textit{Fixed costs:} R\&D of \$14,232m a year \cite[p.~55]{AR} and a manufacturing network of 31 sites \cite[p.~33]{AR}. \textit{Differentiation:} products protected by patents and leading positions until exclusivity ends \cite[pp.~9, 56]{AR}. \textit{Growth:} the sector is growing, with spending up 7.1\% a year to 2029 and sales up 10.0\% in 2025 \cite[p.~6]{AR}, which eases rivalry. More than 100 Phase III studies are ongoing \cite[p.~3]{AR}. & \bbad{Strong} \\
\textbf{New entrants} & \textit{Capital and time:} developing a medicine is estimated at \$2.6bn (in 2013 dollars) \cite{MULLARD}, and R\&D takes 5--15 years and needs regulatory approval \cite[p.~9]{AR}. \textit{Scale:} sales in more than 125 countries with 47,400 commercial employees \cite[p.~32]{AR}, and more than 1,000 collaborations \cite[p.~37]{AR}. \textit{Biotechs:} the Group strengthens its pipeline through acquisitions, licensing and collaborations \cite[p.~48]{AR}, for example EsoBiotec \cite[p.~37]{AR}. \textit{Condition:} the AR warns that limits on intellectual-property protection would lead to ``accelerated and intensified competition'' \cite[p.~49]{AR}. \textit{Exit barriers:} the Group operates specialised manufacturing sites, and the AR records the closure of one tablet plant (Bangalore, May 2025) \cite[p.~33]{AR}. & \bgood{Weak}\newline{\scriptsize conditional on IP} \\
\textbf{Suppliers} & \textit{Concentration:} most API manufacturing is externally sourced, complemented by growing internal capability \cite[p.~33]{AR}. India accounts for 48\% of active API Drug Master Files filed with the US FDA and China for 18\% in 2024; the files reflect manufacturer location, not volume \cite{USP}. \textit{Substitutes:} some inputs have no easy alternative, such as horseshoe-crab blood reagents until a synthetic alternative was adopted in 2025 \cite[p.~45]{AR}. \textit{Mitigation:} dual sourcing, multiple suppliers and stock monitoring \cite[p.~48]{AR}. & \bwarn{Moderate} \\
\textbf{Buyers} & \textit{Concentration:} payers are mainly insurers and national budgets. Payer consolidation and restrictive reimbursement drive significant pricing pressure, and rebates apply \cite[p.~32]{AR}. In the US, three pharmacy benefit managers (CVS Caremark, Express Scripts and Optum Rx) processed about 80\% of equivalent prescription claims in 2025 \cite{FEIN}. Prices are negotiated under the IRA \cite[p.~33]{AR} and follow a four-tier policy based on national income \cite[p.~41]{AR}. Products are prescribed by physicians, who do not pay for them \cite[p.~32]{AR}. Pricing pressure is a Principal Risk \cite[p.~48]{AR}. & \bbad{Strong} \\
\textbf{Substitutes} & \textit{Cost:} when exclusivity ends, generics enter at lower prices. Brilinta sales fell 38\% to \$823m after generic entry in the US and Europe \cite[pp.~18, 56]{AR}, and Symbicort eroded in the EU and Japan \cite[p.~20]{AR}. Farxiga still grew 10\% in 2025 despite generic competition in some markets \cite[p.~56]{AR}. IQVIA estimates more than \$90bn in sales set to lose exclusivity in 2025--2029 \cite{IQVIA}. Generics are treated here as substitutes, because they replace the branded product with the same active ingredient at a lower price \cite[p.~24]{PLAN}. & \bbad{Strong}\newline{\scriptsize over time, by product} \\
\end{xltabular}

\begin{analysis}
For companies already inside, the sector is attractive: it is very hard to enter. But each product only earns well for a limited time. While a medicine is protected, the payers have the upper hand; once the patent expires, generics take over. So the real game is to keep renewing the portfolio. Two ratings depend on where the boundary is drawn: if generics were counted as new entrants, that force would be stronger, and the low threat of entry only holds while patents stay protected \cite[p.~49]{AR}.
\end{analysis}

\subsection{Inside the company: value chain}
The value chain splits what AstraZeneca does into primary activities, which make and deliver the medicines, and support activities, which make the rest possible. The course material also links these activities to the cost structure and profit \cite[p.~25]{PLAN}, and asks what the company does well and less well, including its people and culture \cite[p.~26]{PLAN}.


\begin{table}[htbp]\centering
\caption{Where each dollar of revenue goes in 2025, as a \% of Total Revenue, calculated from the income statement \cite[p.~125]{AR}. Gross margin is 81.9\%.}
\label{tab:custos}
\begin{tabular}{@{}lrB@{}}
\toprule
\textbf{Line} & \textbf{\% of revenue} & \\
\midrule
Selling and administrative & 33.9 & \barra{1.000} \\
R\&D                       & 24.2 & \barra{0.714} \\
Cost of sales              & 18.1 & \barra{0.534} \\
Distribution               &  1.0 & \barra{0.029} \\
\midrule
\textbf{Operating margin}  & \textbf{23.4} & \barrab{0.690} \\
\bottomrule
\end{tabular}
\end{table}
\needspace{16\baselineskip}
\begin{xltabular}{\linewidth}{@{}N{2.45cm}YYp{2.05cm}@{}}
\caption{Value chain analysis. Ratings are not given in the Annual Report.}\label{tab:cadeia-valor}\\*
\toprule
\textbf{Activity} & \textbf{Evidence} & \textbf{Weak points} & \textbf{Rating} \\*
\midrule
\endfirsthead
\toprule
\textbf{Activity} & \textbf{Evidence} & \textbf{Weak points} & \textbf{Rating} \\*
\midrule
\endhead
\bottomrule
\endfoot
\multicolumn{4}{@{}l@{}}{\textbf{Primary activities}} \\*
\midrule
Inbound logistics & Most API manufacturing is externally sourced, complemented by growing internal capability, and a supply chain finance programme supports external suppliers \cite[p.~33]{AR}. & Drug shortages reached record levels in 2025 \cite[p.~33]{AR}; failure to maintain supply is a Principal Risk \cite[p.~48]{AR}. & \bwarn{Neutral / risk} \\
Operations & 31 operations sites in 15 countries (excluding clinical supply sites) and more than 99\% supply performance \cite[pp.~4, 33]{AR}. 63 health-authority GMP and GDP inspections with zero critical findings \cite[p.~31]{AR}. The CEO review cites zero patient-level recalls and zero critical observations in 42 external inspections \cite[p.~4]{AR}. & 10 product recalls \cite[p.~31]{AR}. The AR does not reconcile the 42 and 63 inspection counts. & \bgood{Strength} \\
Outbound logistics & Distribution expense was \$579m, 1.0\% of revenue (calculated) \cite[p.~125]{AR}. 217 on-time market launches are reported under Operations \cite[pp.~27, 33]{AR}. & The AR gives no delivery-performance data. & \bneut{Neutral} \\
Marketing and sales & 47,400 commercial employees \cite[p.~27]{AR}; 16 medicines above \$1bn \cite[p.~4]{AR}; sales through wholly-owned marketing companies, distributors and local offices \cite[p.~32]{AR}. & Five products are 51\% of revenue (calculated: \$30,045m ÷ \$58,739m) and Farxiga alone is 14.5\% \cite[p.~56]{AR}. The US is 43\% of revenue, with price pressure \cite[pp.~2, 32]{AR}. & \bwarn{Strength} \newline{\scriptsize concentration risk} \\
After-sales & Pharmacovigilance across the product life cycle \cite[p.~31]{AR}; patient assistance programmes \cite[p.~41]{AR}; more than 200 partnerships with health systems in 50 countries \cite[p.~4]{AR}. & 10 product recalls \cite[p.~31]{AR}. & \bgood{Strength} \\
\midrule
\multicolumn{4}{@{}l@{}}{\textbf{Support activities}} \\*
\midrule
R\&D and technology & R\&D expense of \$14,232m, 24.2\% of revenue (calculated) \cite[pp.~28, 125]{AR}; 197 pipeline projects \cite[p.~30]{AR}; over 90\% of small-molecule discovery AI-enabled \cite[p.~36]{AR}. & 20 projects discontinued in 2025, 11 for safety or efficacy \cite[p.~30]{AR}; Fasenra's Phase III COPD trial missed its primary endpoint \cite[p.~20]{AR}. & \bgood{Strength} \\
Human resources and culture & 86\% say AstraZeneca is a great place to work (Pulse survey) \cite[pp.~5, 38]{AR}; 89\% report learning or development \cite[p.~39]{AR}; over 50,000 employees took part in the AI programme \cite[pp.~5, 39]{AR}. The Group's five values are in \cite[p.~8]{AR}. & 4,441 Code of Ethics breaches recorded and 434 people removed from role \cite[p.~34]{AR}. Turnover is not reported. & \bgood{Strength} \\
Financial management & Operating cash flow of \$14,575m (+23\%) \cite[p.~10]{AR}. Capital priorities: invest in the business and pipeline, keep an investment-grade rating, pursue business development, and pay a progressive dividend \cite[p.~10]{AR}. & Net debt of \$23,374m, 1.6 times operating cash flow (calculated) \cite[p.~61]{AR}. & \bwarn{Strength} \newline{\scriptsize leverage to monitor} \\
Physical resources & A \$4.5bn plant in Virginia and a \$2bn expansion in Maryland \cite[pp.~4, 33]{AR}. Scope 1 and 2 emissions are 88.1\% below 2015 \cite[p.~42]{AR}. & Against a target of $-$98\% by 2026 \cite[p.~42]{AR}. The Bangalore tablet plant closed in May 2025 \cite[p.~33]{AR}. & \bgood{Strength} \\
Procurement & Dual sourcing, multiple suppliers and stock monitoring \cite[p.~48]{AR}. & The AR describes mitigation, not performance. & \bneut{Neutral} \\
Administration and risk & Each SET function assesses risk quarterly, feeding the Group Risk Report for the Board, the Audit Committee and the SET \cite[p.~47]{AR}. The PAAGR restructuring programme has realised annual run-rate pre-tax benefits of \$1.9bn \cite[p.~59]{AR}. & All significant operating decisions are taken at SET level \cite[p.~147]{AR} & \bwarn{Neutral / risk} \\
\midrule
\multicolumn{4}{@{}l@{}}{\textbf{Margin}} \\*
\midrule
Cost structure and margin & 2025, \% of Total Revenue (calculated from \cite[p.~125]{AR}): cost of sales 18.1\%, distribution 1.0\%, R\&D 24.2\%, selling, general and administrative 33.9\%. Gross margin 81.9\% and operating margin 23.4\%. & --- & \bgood{Strength} \\
\end{xltabular}

\begin{analysis}
Most of the chain works well: supply is reliable, research spending is high and the business generates plenty of cash. The weak spots sit at the edges: raw materials depend on outside suppliers, sales depend on a handful of products, and some research projects fail. The course material asks us to compare with competitors \cite[p.~24]{PLAN}, but the only peer comparison in the Annual Report is shareholder return over ten years: 307\%, against 197\% for global pharma peers and 65\% for European peers \cite[p.~90]{AR}. For everything else the Group can only be compared with its own record and targets.
\end{analysis}

\subsection{Which resources give an edge: VRIO}
VRIO asks four questions about each resource: does it create value, is it rare in the sector, is it hard or costly to copy, and is the company organised to use it? Only a resource that passes all four counts as a competitive advantage \cite[p.~26]{PLAN}. Because the Annual Report does not compare AstraZeneca with its rivals, the answers below are an assessment, not a company statement.

\begin{xltabular}{\linewidth}{@{}N{2.4cm}Yccccp{2.5cm}@{}}
\caption{VRIO analysis: Value, Rarity, Imitability, Organisation. The result follows the course rule \cite[p.~26]{PLAN}.}\label{tab:vrio}\\*
\toprule
\textbf{Resource} & \textbf{Why} & \textbf{V} & \textbf{R} & \textbf{I} & \textbf{O} & \textbf{Result} \\*
\midrule
\endfirsthead
\toprule
\textbf{Resource} & \textbf{Why} & \textbf{V} & \textbf{R} & \textbf{I} & \textbf{O} & \textbf{Result} \\*
\midrule
\endhead
\bottomrule
\endfoot
\textbf{Pipeline and R\&D} & 197 projects; 20 new medicines in late stages or under review; 43 major approvals in 2025 in the US, EU, China and Japan \cite[pp.~28, 30]{AR} & Y & Y & Y & Y & \bgood{Advantage} \\
\textbf{Leadership in oncology} & \$25.6bn (+15\%) and five medicines above \$1bn; the leadership claims are the Group's own \cite[pp.~4, 12--14]{AR} & Y & Y & Y & Y & \bgood{Advantage} \\
\textbf{Patents on current products} & The law blocks copies while they last; Farxiga loses US exclusivity in 2026 \cite[pp.~9, 33]{AR} & Y & Y & Y & Y & \bwarn{Until expiry} \\
\textbf{Operational excellence} & >99\% supply performance and 217 on-time launches \cite[pp.~4, 27]{AR} & Y & Y & N & Y & \bneut{Parity} \newline{\scriptsize can be copied} \\
\textbf{Global presence} & Sales in over 125 countries; Emerging Markets are 26\% of revenue \cite[pp.~2, 32--33]{AR} & Y & N & N & Y & \bneut{Parity} \newline{\scriptsize rivals have it too} \\
\textbf{AI and data} & 90\%+ of small-molecule discovery uses AI; Enterprise AI unit created in December 2025; scaling AI is listed as a Principal Risk \cite[pp.~36, 48]{AR} & Y & N & N & Y & \bneut{Parity} \newline{\scriptsize sector-wide} \\
\textbf{Partnership network} & More than 1,000 collaborations worldwide \cite[p.~37]{AR} & Y & N & N & Y & \bneut{Parity} \newline{\scriptsize rivals have it too} \\
\end{xltabular}

\begin{analysis}
AstraZeneca's real edge is its \textbf{ability to innovate}: the pipeline, R\&D and its position in oncology. Patents also protect it, but only until they expire. Operations, global reach, AI and partnerships are valuable, yet competitors have or can build them too, so they do not set the company apart. This matches what Porter showed: success depends on renewing the portfolio.
\end{analysis}

\subsection{Putting it together: the SWOT}
The SWOT gathers the most important points from the four tools. From the PESTEL, the factors with high or medium impact are kept; from the other tools, the points that matter most or that come up more than once. Under each point, a grey label shows which tool it came from, so the origin of the information is always visible. Strengths whose label is marked parity passed the value test but failed the rarity or the imitability test of the VRIO: they are needed to compete, but they do not set the company apart.

\begin{xltabular}{\linewidth}{@{}lp{3.6cm}Y@{}}
\caption{SWOT matrix of AstraZeneca, with the source tool of each point.}\label{tab:swot}\\*
\toprule
\textbf{Code} & \textbf{Factor} & \textbf{Evidence and origin} \\*
\midrule
\endfirsthead
\toprule
\textbf{Code} & \textbf{Factor} & \textbf{Evidence and origin} \\*
\midrule
\endhead
\bottomrule
\endfoot
\multicolumn{3}{@{}l@{}}{\textbf{Strengths --- inside the company, today}} \\*
\midrule
S1 & \textbf{Strong pipeline and R\&D} & \$14.2bn invested (24\% of revenue), 197 projects, 20 new medicines in late stages \cite[pp.~28, 30, 125]{AR} \newline \textit{\small VRIO · Pipeline\quad Value chain · R\&D} \\
S2 & \textbf{Leadership in oncology} & 44\% of revenue, growing 15\% \cite[pp.~2, 4, 13]{AR} \newline \textit{\small VRIO · Oncology} \\
S3 & \textbf{Global reach, broad portfolio} & Sales in 125+ countries; 16 medicines above \$1bn \cite[pp.~4, 32]{AR} \newline \textit{\small VRIO · Global presence: parity\quad Value chain · Sales} \\
S4 & \textbf{Reliable operations} & >99\% supply, 217 on-time launches, zero critical findings in 63 inspections \cite[pp.~4, 27, 31]{AR} \newline \textit{\small VRIO · Operations: parity\quad Value chain · Operations} \\
S5 & \textbf{Financial capacity and margin} & Gross margin of 81.9\% and operating margin of 23.4\%; operating cash flow of \$14.6bn (+23\%), against net debt of \$23.4bn \cite[pp.~10, 61, 125]{AR} \newline \textit{\small Value chain · Finance} \\
S6 & \textbf{Advanced use of AI} & 90\%+ of small-molecule discovery uses AI; dedicated AI unit \cite[pp.~5, 36]{AR} \newline \textit{\small VRIO · AI: parity\quad Value chain · R\&D} \\
S7 & \textbf{Proven ability to restructure} & The PAAGR programme has realised annual run-rate pre-tax benefits, before reinvestment, of \$1.9bn \cite[p.~59]{AR} \newline \textit{\small Value chain · Administration} \\
\midrule
\multicolumn{3}{@{}l@{}}{\textbf{Weaknesses --- inside the company, today}} \\*
\midrule
W1 & \textbf{Too much revenue from five products} & Farxiga, Tagrisso, Imfinzi, Ultomiris and Calquence make \$30,045m: 51\% of Total Revenue, or 54\% of Product Sales on a comparable basis (calculated from \cite[pp.~56, 125]{AR}) \newline \textit{\small Value chain · Sales} \\
W2 & \textbf{Heavy reliance on the US} & 43\% of revenue, heading to about 50\% by 2030 \cite[pp.~2, 4, 32]{AR} \newline \textit{\small Value chain · Sales} \\
W3 & \textbf{Older products already losing ground} & Brilinta $-$38\% after generics; Symbicort eroding in the EU and Japan \cite[pp.~18, 20, 56]{AR} \newline \textit{\small Porter · Substitutes\quad Value chain · Sales} \\
W4 & \textbf{Slow or negative growth in some areas} & Cardiovascular, renal \& metabolism +3\% (\$12.8bn, 22\% of revenue); vaccines $-$13\% \cite[pp.~4, 17, 56]{AR} \newline \textit{\small Value chain · Sales} \\
W5 & \textbf{Behind on its 2026 climate target} & $-$88.1\% emissions against a $-$98\% target; climate is not a Principal Risk in the AR \cite[pp.~42, 47]{AR} \newline \textit{\small Value chain · Facilities} \\
W6 & \textbf{Setbacks in the pipeline} & 20 projects stopped in 2025, 11 for safety or efficacy; Fasenra failed a Phase III COPD trial \cite[pp.~20, 30]{AR} \newline \textit{\small Value chain · R\&D} \\
W7 & \textbf{Operating decisions concentrated in the SET} & All significant operating decisions are taken at SET level \cite[p.~147]{AR}, which slows the local response that the strategy depends on \newline \textit{\small Value chain · Administration} \\
\midrule
\multicolumn{3}{@{}l@{}}{\textbf{Opportunities --- outside the company, future}} \\*
\midrule
O1 & \textbf{Rising demand} & Ageing and chronic disease \cite[pp.~6--7]{AR} \newline \textit{\small PESTEL · Social} \\
O2 & \textbf{Growing health spending} & +7.1\% a year to 2029; pharma sales +10\% in 2025 \cite[p.~6]{AR} \newline \textit{\small PESTEL · Economic} \\
O3 & \textbf{Emerging markets} & \$15.3bn (+12\%), +19\% outside China \cite[pp.~4, 33]{AR} \newline \textit{\small PESTEL · Economic} \\
O4 & \textbf{AI and new technology} & \$4--7bn a year for the sector; access through 1,000+ collaborations \cite[pp.~7, 37]{AR} \newline \textit{\small PESTEL · Technological} \\
O5 & \textbf{Large unmet needs} & Hypertension, cancer, and rare diseases where fewer than 10\% have a treatment \cite[pp.~13, 17, 23]{AR} \newline \textit{\small PESTEL · Social} \\
O6 & \textbf{US deal that rewards local production} & Three-year tariff exemption, alongside prices aligned with other wealthy countries \cite[pp.~4, 32--33]{AR} \newline \textit{\small PESTEL · Political} \\
O7 & \textbf{The sector is hard to enter} & Developing a medicine costs billions and takes 5 to 15 years, on top of the regulatory barrier, so the threat of new entrants is weak and value stays with those already inside \cite{MULLARD}, \cite[p.~9]{AR} \newline \textit{\small Porter · New entrants} \\
\midrule
\multicolumn{3}{@{}l@{}}{\textbf{Threats --- outside the company, future}} \\*
\midrule
T1 & \textbf{Price pressure} & Tight budgets, IRA negotiation (Farxiga 2026, Calquence 2027), consolidated payers \cite[pp.~7, 32--33]{AR}, \cite{FEIN} \newline \textit{\small PESTEL · Economic, Legal\quad Porter · Buyers} \\
T2 & \textbf{Patents running out} & Farxiga (14.5\% of revenue) loses US exclusivity in 2026; the AR expects this to be manageable. \$90bn+ of sector sales at risk to 2029 \cite[pp.~33, 56]{AR}, \cite{IQVIA} \newline \textit{\small Porter · Substitutes\quad VRIO · Patents} \\
T3 & \textbf{Geopolitics and tariffs} & US--China rivalry and trade controls \cite[pp.~6, 49]{AR} \newline \textit{\small PESTEL · Political} \\
T4 & \textbf{Intense rivalry} & Large rivals of similar size \cite{P50}, \cite[p.~3]{AR} \newline \textit{\small Porter · Rivalry} \\
T5 & \textbf{Weaker patents, new EU rules, approval delays} & \cite[pp.~48, 49]{AR}, \cite{EU} \newline \textit{\small PESTEL · Legal} \\
T6 & \textbf{Supply and climate disruption} & Record drug shortages, outsourced APIs concentrated in India and China, extreme weather \cite[pp.~7, 33, 44]{AR}, \cite{USP} \newline \textit{\small Porter · Suppliers\quad PESTEL · Environmental} \\
\end{xltabular}

\subsection{What matters most}
Not all 27 points weigh the same. They are ranked by their weight in revenue, how soon they hit, whether the Annual Report treats them as a priority or Principal Risk, and whether they can be measured.

\medskip
\noindent\textbf{Inside the company}
\begin{enumerate}[topsep=4pt,itemsep=2pt]
  \item 51\% of revenue from five products, 14.5\% from Farxiga alone; the other side of the strong oncology portfolio \cite[pp.~2, 56]{AR}
  \item 9 of 20 new medicines promised for 2030 already delivered, backed by \$14.2bn of R\&D. The AR calls R\&D success crucial \cite[pp.~10, 28, 52]{AR}
  \item \$14.6bn of operating cash flow and >99\% supply fund and deliver the \$50bn US plan, against \$23.4bn of net debt \cite[pp.~4, 10, 61]{AR}
\end{enumerate}

\medskip
\noindent\textbf{Outside the company}
\begin{enumerate}[topsep=4pt,itemsep=2pt]
  \item 2026 Farxiga loses US exclusivity and gets a negotiated price in the same year; Calquence follows in 2027 \cite[pp.~33, 56]{AR}
  \item 43\% $\rightarrow$ 50\% of revenue in the US by 2030 is the company's own exposure (W2); what sits outside is the policy acting on it --- Medicare-negotiated prices and payers concentrated in three pharmacy benefit managers that handle about 80\% of US prescriptions \cite[pp.~2, 4, 32]{AR}, \cite{FEIN}
  \item +7.1\% a year in health spending and a growing disease burden: the counterweight to the two threats above \cite[pp.~6, 7, 13, 17]{AR}
\end{enumerate}

\begin{analysis}
\textbf{The link between the two}\par
Everything comes down to one race: launching new medicines faster than old ones lose protection. Brilinta has already lost 38\% of its sales and Farxiga, the biggest product, is about to lose US exclusivity \cite[pp.~33, 56]{AR}. The market keeps growing in value, but each product earns for a shorter time and at a lower price.
\end{analysis}

\subsection{Crossing the factors}
The course material suggests four questions to link the inside with the outside \cite[pp.~27--28]{PLAN}. Each cell answers one of them.

\begin{table}[htbp]\centering
\caption{The four crossings between the inside and the outside. The cells are not given in the Annual Report.}
\label{tab:cruzamentos}
\begin{tabularx}{\linewidth}{@{}l>{\hsize=.5\hsize}Y>{\hsize=1.5\hsize}Y@{}}
\toprule
& \textbf{Crossing} & \textbf{What it means for AstraZeneca} \\
\midrule
\multirow{2}{*}{\rotatebox[origin=c]{90}{\textbf{\footnotesize Strengths}}}
 & \textbf{Grow} (SO) & Use the pipeline and oncology (S1, S2) to meet growing demand, rising health spending and unmet needs (O1, O2, O5), and AI (S6) to speed up research (O4). Entry barriers keep those gains with the companies already inside (O7). \\
\addlinespace
 & \textbf{Defend} (ST) & Use global reach (S3) against tariffs and geopolitics (T3), renew the portfolio (S1) before patents expire (T2), and use cash and the proven ability to restructure (S5, S7) to fund US plants and acquisitions against tariffs (T3). \\
\midrule
\multirow{2}{*}{\rotatebox[origin=c]{90}{\textbf{\footnotesize Weaknesses}}}
 & \textbf{Fix} (WO) & Reduce dependence on five products (W1) by entering areas of large need such as hypertension (O5), where baxdrostat has two positive Phase III results \cite[p.~17]{AR}. Emerging markets (O3) add growth, but with the same lead products \cite[p.~56]{AR}. \\
\addlinespace
 & \textbf{Avoid} (WT) & Concentration and fading products (W1, W3) combined with price pressure and expiring patents (T1, T2) are the most dangerous mix, and so is US reliance (W2) under price and policy pressure (T1, T3). Decisions concentrated in the SET (W7) slow the local response that price pressure (T1) demands. \\
\bottomrule
\end{tabularx}
\end{table}

\subsection{Proposed strategy}
The course material asks for several alternatives to be weighed before one is chosen \cite[pp.~3, 7, 29]{PLAN}. It does not offer a list of generic strategies, and the Annual Report does not reject any option, so the three levers the report itself describes are compared below.

\begin{table}[htbp]\centering
\caption{Alternatives considered before choosing the strategy.}
\label{tab:alternativas}
\begin{tabularx}{\linewidth}{@{}lp{3.5cm}Y@{}}
\toprule
\textbf{Role} & \textbf{Lever} & \textbf{Evidence and what it answers} \\
\midrule
Core & Organic growth through R\&D & New medicines and new uses of existing ones; 12 first life-cycle approvals in 2025 \cite[pp.~4, 52]{AR}. Answers S1, S2 $\rightarrow$ O1, O5, O7, T2. \\
\addlinespace
Complement & Acquisitions, licences and partnerships & For example EsoBiotec, up to \$1bn \cite[pp.~37, 48]{AR}. Answers T2, W1, W3. \\
\addlinespace
Kept as policy & Paying shareholders & Progressive dividend of \$3.20 a share \cite[pp.~3, 10]{AR}. Important, but it does not fix any weakness or threat. \\
\bottomrule
\end{tabularx}
\end{table}
The proposed strategy is \textbf{growth through innovation, with faster renewal of the portfolio, better geographic balance and secure supply}. It draws on all four crossings in five actions. The growth itself rests on S1 and S2, the two resources the VRIO confirms as advantages that do not expire; S3, S4 and S6 are conditions for carrying the strategy out, not sources of advantage in themselves.

\begin{enumerate}[topsep=4pt,itemsep=3pt]
  \item \textbf{Focus R\&D where the company leads or wants to lead} Oncology, cardiovascular, renal and metabolic, respiratory and immunology, and rare diseases, using AI to cut development time. Late-stage medicines such as baxdrostat also address slow growth in cardiovascular \cite[p.~17]{AR}. \textit{\small S1, S2, S6 $\rightarrow$ O1, O4, O5 · W4}
  \item \textbf{Replace the revenue that patents will take away} With new medicines and new indications, topped up with acquisitions, licences and partnerships \cite[pp.~37, 48]{AR}. \textit{\small S1, S5 $\rightarrow$ T2, W1, W3}
  \item \textbf{Grow in emerging markets and produce locally where it pays} Emerging markets bring \$15.3bn and grow 12\% \cite[pp.~4, 33]{AR}; local plants in the US, China and Singapore \cite[pp.~4, 5, 32, 37]{AR}. This pulls against the US share rising to about 50\% by 2030 \cite[p.~4]{AR}. \textit{\small S3 $\rightarrow$ O3, O6, T3}
  \item \textbf{Show payers what the medicines are worth, and let regions adapt} Keep the four-tier pricing policy and the outcome-based agreements \cite[pp.~41, 48]{AR}, but give regional commercial teams room to adapt access and pricing within limits set centrally. All significant operating decisions are taken at SET level today \cite[p.~147]{AR}; delegating this one brings the decision closer to the payer that has to be convinced, and the PAAGR shows the Group can reorganise at that scale \cite[p.~59]{AR}. \textit{\small S7 $\rightarrow$ T1 · W7}
  \item \textbf{Secure supply and deliver the climate plan} Keep supply above 99\% \cite[p.~4]{AR}, add capacity \cite[p.~33]{AR} and close the gap to the $-$98\% emissions target \cite[p.~42]{AR}. \textit{\small S4 $\rightarrow$ T6, W5}
\end{enumerate}

\begin{figure}[htbp]\centering
\begin{tikzpicture}[node distance=6mm and 6mm]
\node[fl] (c1) {Pipeline, R\&D and AI};
\node[fl, right=of c1] (c2) {New medicines, acquisitions and licences};
\node[fl, right=of c2] (c3) {Replace revenue losing protection};
\node[fle, right=of c3] (c4) {Ambition 2030\\{\color{azmuted}$\geq$20 new medicines, \$80bn}};
\node[fl, below=of c1] (d1) {Global presence and local plants};
\node[fl, right=of d1] (d2) {Emerging markets, fewer tariff risks};
\node[fl, right=of d2] (d3) {Value for payers, secure supply, climate plan};
\node[fle, right=of d3] (d4) {Sustainable growth};
\draw[se] (c1)--(c2); \draw[se] (c2)--(c3); \draw[se] (c3)--(c4);
\draw[se] (d1)--(d2); \draw[se] (d2)--(d3); \draw[se] (d3)--(d4);
\end{tikzpicture}
\caption{How the five actions lead to the Ambition 2030. The 2030 targets are the Group's \cite[p.~10]{AR}.}
\label{fig:acoes}
\end{figure}
Two weaknesses stay without a dedicated action: US reliance (W2) and pipeline setbacks (W6), which are treated here as risks to watch. Rivalry and patent rules (T4, T5) are outside the company's control, and the strategy answers them indirectly by renewing the portfolio.

How does this compare with what AstraZeneca already says? Its declared strategy, \textit{Growth Through Innovation}, has three priorities \cite[pp.~2, 10, 11]{AR}, and the Annual Report lists similar mitigations for its main risks \cite[p.~48]{AR}. Actions 1, 2 and 5 stay close to that strategy. Actions 3 and 4 depart from it: Action 3 puts more weight on emerging markets than the company does, and Action 4 proposes delegating pricing decisions that are taken centrally today. What this adds is the explicit path from the SWOT to the crossings to action, and the choice between levers.

\begin{xltabular}{\linewidth}{@{}N{2.5cm}Yp{2.2cm}Y@{}}
\caption{Indicators to follow the strategy. Targets are declared by the Group; the choice of indicators is not. The Annual Report sets no target for revenue concentration or for the slow-growth areas. The scorecard cash-flow measure is operating cash flow less capital expenditure, plus proceeds from disposals of intangible assets \cite[p.~91]{AR}, which is why it is lower than the \$14.6bn in S5.}\label{tab:indicadores}\\*
\toprule
\textbf{Goal} & \textbf{Indicator} & \textbf{Today} & \textbf{Declared target} \\*
\midrule
\endfirsthead
\toprule
\textbf{Goal} & \textbf{Indicator} & \textbf{Today} & \textbf{Declared target} \\*
\midrule
\endhead
\bottomrule
\endfoot
Renew the portfolio & New medicines approved since Oct 2022 & 9 & $\geq$ 20 by 2030 \cite[p.~10]{AR} \\
Grow & Total Revenue & \$58.7bn (2025) & \$80bn by 2030 \cite[p.~10]{AR} \\
R\&D productivity & Regulatory events per year & 97 (2025) & 2025 target 54, with 69 counted; 2026 not disclosed \cite[pp.~91, 102]{AR} \\
R\&D productivity & Pipeline progression events & 38 (2025) & 2025 target 29, with 36 counted \cite[pp.~11, 91]{AR} \\
Financial capacity & Cash flow (Group scorecard measure) & \$11.9bn (2025) & 2025 target \$11.3bn \cite[p.~91]{AR} \\
Sustainability & Scope 1 and 2 emissions vs 2015 & $-$88.1\% & $-$98\% by 2026 \cite[p.~42]{AR} \\
Sustainability & Scope 3 emissions vs 2019 & --- & $-$50\% by 2030, $-$90\% by 2045; net zero by 2045 \cite[p.~42]{AR} \\
\end{xltabular}
